\section{Aufbau des Frameworks}
In diesem Kapitel wird zunächst der technische Aufbau des Frameworks beschrieben.  
Darauf aufbauend wird die Softwarearchitektur erläutert,
wobei das Zusammenspiel der zentralen Module im Rahmen des Model-View-Controller-Paradigmas dargestellt wird.

\subsection{Technischer Aufbau}
Das entwickelte Framework ist als eigenständiges Python-Paket organisiert und wurde mithilfe des Packaging-Werkzeugs poetry erstellt.
Diese Struktur ermöglicht eine modulare Organisation des Quellcodes,
die saubere Trennung von Abhängigkeiten sowie eine einfache Verteilung und Wiederverwendbarkeit des Frameworks \doublecite[1]{python_packages}.
Die Verzeichnisstruktur folgt dem üblichen Aufbau für Python-Pakete mit einem root directory, den entsprechenden Python Modulen,
einer pyproject.toml, einer README.md sowie einem Verzeichnis für Tests \doublecite[22]{python_packages}:

\vspace{1em}
\begin{figure}[H]
    \dirtree{%
    .1 visualkinematics\_v2.
    .2 poetry.lock.
    .2 pyproject.toml.
    .2 README.md.
    .2 assets.
    .3 custom-package-sets.
    .3 ros-package-sets.
    .2 src.
    .3 examples.
    .4 industrial\_robot\_example.ipynb.
    .4 quader\_example.ipynb.
    .4 robot\_example.ipynb.
    .4 rotations\_example.ipynb.
    .3 visualkinematics\_v2.
    .4 \_\_init\_\_.py.
    .4 manager.py.
    .4 manipulator.py.
    .4 show.py.
    .4 util.py.
    .4 \_package\_data.
    .2 tests.
    }
\caption{Verzeichnisstruktur des Pakets visualkinematics\_v2.}
\label{fig:dir_structure}
\end{figure}
\vspace{1em}


Durch diese Paketstruktur kann das Framework problemlos in andere Projekte integriert, installiert oder erweitert werden.
Vorgefertigte Beispiel-Notebooks befinden sich im Ordner examples und veranschaulichen die Funktionsweise des Frameworks.
Darunter befinden sich die Notebooks quader\_example.ipynb, robot\_example.ipynb sowie rotations\_example.ipynb,
die bereits im Rahmen der Projektarbeit entwickelt worden sind und die Funktionsweise von Transformationen von einfachen Körpern
sowie von mobilen Robotern in Form von Übungsaufgaben demonstrieren.
Das Beispiel-Notebook industrial\_robot\_example.ipynb stellt die Nutzung des Frameworks zur Simulation von Industrierobotern dar.
Alle Beispiel-Notebooks greifen direkt auf den Quellcode des Pakets zu.

Der Quellcode befindet sich im Verzeichnis visualkinematics\_v2/src/visual\-ki\-ne\-ma\-tics\_v2.
Dieser ist in einzelne Module gegliedert, die jeweils einer bestimmten Verantwortlichkeit zugeordnet sind. 
Ein Überblick über die zentralen Module und deren Beziehungen ist in Abbildung \ref{fig:package_diagramm} (siehe Anhang) in Form eines Paketdiagramms dargestellt. 
Es wurde darauf geachtet, dass zwischen den Modulen keine bidirektionalen oder zyklischen Abhängigkeiten bestehen.
Eine ausführliche Dokumentation des Quellcodes liegt der Arbeit als HTML-Datei im Verzeichnis docs bei.

\subsection{Softwarearchitektur Model-View-Controller}
Der Quellcode des Pakets verteilt sich auf insgesamt vier Module: util.py, manipulator.py, manager.py und show.py.
Das Zusammenspiel dieser Module folgt der Softwarearchitektur des Model-View-Controller (MVC) Paradigmas,
welches sich für grafische Echtzeitanwendungen eignet.
Ziel dieser Trennung ist eine modulare Struktur, die eine klare Abgrenzung zwischen Daten (Model),
Präsentation (View) und Interaktion/Steuerung (Controller) ermöglicht.
Diese Struktur erleichtert die Erweiterbarkeit und Wartbarkeit des Frameworks \doublecite[1000]{Encyclopedia_of_Computer_Graphics_and_Games}.




\subsubsection{Model-Komponente}
Das Modul util.py stellt eine zustandslose Sammlung fachspezifischer Kern- und Hilfsfunktionen dar und lässt sich am ehesten der Model-Komponente zuordnen.
Es wird von allen anderen Modulen importiert und genutzt.
Die Model-Komponente wird vor allem durch die Module manipulator.py und manager.py realisiert.  
Sie bildet die strukturelle und funktionale Grundlage des Frameworks und kapselt die kinematische Struktur des Roboters sowie dessen aktuellen Zustand.

Das Modul manipulator.py enthält zentrale Klassen wie:

\begin{itemize}
    \item \texttt{Manipulator} – verwaltet den Gesamtaufbau des Roboters, einschließlich aller Glieder und Gelenke.
    \item \texttt{Link} – repräsentiert ein einzelnes physisches Glied der kinematischen Kette und enthält Informationen zur Geometrie und lokalen Transformation.
    \item \texttt{Joint} – modelliert ein Gelenk zwischen zwei Gliedern inklusive Bewegungsbereich und aktuellem Winkel.
    \item \texttt{Kinematic\_Chain\_Element} – die Oberklasse von Link und Joint und verwaltet die Children und den Parent die sich aus der kinematischen Kette ergeben.
    \item \texttt{DHKinematicModel} – kapselt die mathematische Berechnung der Vorwärts- und Rückwärtstransformation auf Basis der Denavit-Hartenberg-Parameter.
\end{itemize}

Innerhalb dieser Klassen werden alle relevanten Informationen zur Beschreibung und Steuerung eines Industrieroboters zusammengeführt.  
Dazu zählen die Hierarchie der Glieder, die Relativtransformationen entlang der kinematischen Kette sowie der aktuelle Gelenkzustand (z.\,B. Positionen und Winkel der Achsen).  
Diese Daten bilden die Grundlage für die spätere grafische Darstellung und Animation.
Die Denavit-Hartenberg-Konvention sowie die darauf basierende Berechnung der Vorwärts- und Rückwärtstransformation wurden in die Klasse DHKinematicModel ausgelagert.
Dadurch wird eine saubere Trennung der Zuständigkeiten erreicht und gleichzeitig die Wartbarkeit sowie die Lesbarkeit des Codes verbessert.

Innerhalb der Model-Komponente wird die Datenhaltung und somit der Zugriff auf den persistenten Speicher in das Modul manager.py entkoppelt.
Hier befinden sich Funktionen zum Laden und Speichern von 3D-Modellen sowie URDF-Dateien.
Darüber hinaus erzeugt das Modul manager.py eine Datei namens config.toml in einem standardisierten Verzeichnis für benutzerspezifische Konfigurationsdaten.
Der genaue Speicherort ist betriebssystemabhängig. Unter Windows befindet sich dieser beispielsweise unter
C:/Users/\textless{}Benutzername\textgreater{}/AppData/Local/visualkinematics\_v2,
während auf Linux-Systemen typischerweise das Verzeichnis
~/.config/visualkinematics\_v2 verwendet wird.
Dieses Vorgehen entspricht den Konventionen der jeweiligen Plattformen und ermöglicht eine konsistente sowie benutzerspezifische Speicherung von Anwendungseinstellungen.
In der Datei config.toml werden die Speicherorte für Ressourcen und Assets definiert. Der Benutzer hat die Möglichkeit,
diese Pfade bei Bedarf manuell anzupassen, um beispielsweise benutzerdefinierte Verzeichnisse zu verwenden.
In den konfigurierten Verzeichnissen können benutzereigene 3D-Modelle und URDF-Beschreibungen abgelegt werden.
Diese stehen dem Framework anschließend zur Verfügung und können direkt geladen werden.
Damit das bereitgestellte Beispiel-Notebook industrial\_robot\_ex\-am\-ple.ipynb ohne weitere Anpassungen („out of the box“) lauffähig ist,
werden mit dem Paket im Verzeichnis \_package\_data bereits einige 3D-Modelle und URDF-Be\-schrei\-bun\-gen mitgeliefert.


\subsubsection{View-Komponente}
Die View ist für die Darstellung der Szene im Jupyter Notebook zuständig und wird durch das Modul show.py dargestellt.
Sie verwendet pythreejs sowie ipywidgets zur visuellen Repräsentation und enthält folgende Klassen und Funktionen:

\begin{itemize}
    \item \texttt{Environment}: erzeugt die 3D-Szene mit Licht, Kamera, Koordinatensystem usw.
    \item \texttt{\_ipython\_display\_()}: bindet die Szene beim Aufrufen von display() als interaktives Widget in das Notebook ein.
    \item \texttt{Inspector\_View}: enthält GUI-Elemente wie Schieberegler und Knöpfe für die Steuerung und Konfiguration eines Roboters.
    \item \texttt{Gizmo\_Controls\_View}: enthält GUI-Elemente für die Translation, Rotation und Skalierung eines beliebigen Körpers.
\end{itemize}

\subsubsection{Controller-Komponente}
Der Controller bildet die Schnittstelle zwischen der Benutzerinteraktion und dem zugrunde liegenden Datenmodell.
Er verknüpft die GUI-Elemente der View mit den Funktionen des Models und steuert das Verhalten des Frameworks in Abhängigkeit von Benutzereingaben.
Dabei kommt das Modul ipyevents zum Einsatz, das eine Ereignisverarbeitung in Jupyter Notebooks ermöglicht.
Der Controller ist wie die View-Komponente ebenfalls im Modul show.py implementiert und wird dort in den folgenden Klassen gekapselt:

\begin{itemize}
    \item \texttt{Inspector\_Controller}: Verbindet die GUI-Elemente von Inspector\_View mit Funktionen und speichert Zustände wie mouse\_down.
    \item \texttt{Gizmo\_Controls\_Controller}: Verbindet die GUI-Elemente der \\Giz\-mo\_Con\-trols\_View mit Funktionen.
\end{itemize}

Zwischen der View und dem Controller sowie zwischen dem Model und dem Controller besteht eine lose Kopplung.
Die Klasse Inspector\_Controller besitzt eine Referenz auf die Klasse Manipulator (Model), wodurch eine einseitige Assoziation zwischen Controller und Model besteht.
Die Klasse Inspector\_View kennt die Manipulator-Instanz ebenfalls, allerdings ausschließlich lesend, um beispielsweise Initialwerte für GUI-Elemente wie Schieberegler setzen zu können.
Eine detaillierte Darstellung der beteiligten Module, Klassen und ihrer Beziehungen befindet sich in Abbildung \ref{fig:class_diagramm} (siehe Anhang) in Form eines Klassendiagramms.

